\documentclass[a4paper,11pt]{article}

\usepackage{latexsym}
\usepackage[empty]{fullpage}
\usepackage{titlesec}
\usepackage{marvosym}
\usepackage[usenames,dvipsnames]{color}
\usepackage{verbatim}
\usepackage{enumitem}
\usepackage[hidelinks]{hyperref}
\usepackage{fancyhdr}
\usepackage[italian]{babel}
\usepackage{tabularx}
\input{glyphtounicode}


%----------FONT OPTIONS----------
% sans-serif
% \usepackage[sfdefault]{FiraSans}
% \usepackage[sfdefault]{roboto}
% \usepackage[sfdefault]{noto-sans}
% \usepackage[default]{sourcesanspro}

% serif
% \usepackage{CormorantGaramond}
% \usepackage{charter}


\pagestyle{fancy}
\fancyhf{} % clear all header and footer fields
\fancyfoot{}
\renewcommand{\headrulewidth}{0pt}
\renewcommand{\footrulewidth}{0pt}

% Adjust margins
\addtolength{\oddsidemargin}{-0.5in}
\addtolength{\evensidemargin}{-0.5in}
\addtolength{\textwidth}{1in}
\addtolength{\topmargin}{-.5in}
\addtolength{\textheight}{1.0in}

\urlstyle{same}

\raggedbottom
\raggedright
\setlength{\tabcolsep}{0in}

% Sections formatting
\titleformat{\section}{
  \vspace{-4pt}\scshape\raggedright\large
}{}{0em}{}[\color{black}\titlerule \vspace{-5pt}]

% Ensure that generate pdf is machine readable/ATS parsable
\pdfgentounicode=1

%-------------------------
% Custom commands
\newcommand{\resumeItem}[1]{
  \item\small{
    {#1 \vspace{-2pt}}
  }
}

\newcommand{\resumeSubheading}[4]{
  \vspace{-2pt}\item
    \begin{tabular*}{0.97\textwidth}[t]{l@{\extracolsep{\fill}}r}
      \textbf{#1} & #2 \\
      \textit{\small#3} & \textit{\small #4} \\
    \end{tabular*}\vspace{-7pt}
}

\newcommand{\resumeSubSubheading}[2]{
    \item
    \begin{tabular*}{0.97\textwidth}{l@{\extracolsep{\fill}}r}
      \textit{\small#1} & \textit{\small #2} \\
    \end{tabular*}\vspace{-7pt}
}

\newcommand{\resumeProjectHeading}[2]{
    \item
    \begin{tabular*}{0.97\textwidth}{l@{\extracolsep{\fill}}r}
      \small#1 & #2 \\
    \end{tabular*}\vspace{-7pt}
}

\newcommand{\resumeSubItem}[1]{\resumeItem{#1}\vspace{-4pt}}

\renewcommand\labelitemii{$\vcenter{\hbox{\tiny$\bullet$}}$}

\newcommand{\resumeSubHeadingListStart}{\begin{itemize}[leftmargin=0.15in, label={}]}
\newcommand{\resumeSubHeadingListEnd}{\end{itemize}}
\newcommand{\resumeItemListStart}{\begin{itemize}}
\newcommand{\resumeItemListEnd}{\end{itemize}\vspace{-5pt}}

%-------------------------------------------
%%%%%%  CV INIZIA QUI  %%%%%%%%%%%%%%%%%%%%%%%%%%%%
%-------------------------------------------


\begin{document}

\begin{center}
    \textbf{\Huge \scshape Ledjo Lleshaj} \\ \vspace{1pt}
    \small (+39) 388-244-1632 $|$ \href{mailto:ledjolleshaj@yahoo.com}{\underline{ledjolleshaj@yahoo.com}} $|$
    \href{https://linkedin.com/in/ledjo-lleshaj/}{\underline{linkedin.com/in/ledjo-lleshaj/}} $|$
    \href{https://github.com/LedjoLleshaj}{\underline{github.com/LedjoLleshaj}}
\end{center}


%-----------PROFILO-----------
\section{Profilo}
\small{Backend Software Engineer con Laurea Magistrale in Informatica, esperienza nello sviluppo di API RESTful, applicazioni web full-stack e sistemi integrati con intelligenza artificiale. Competente in Python (Django, Flask, FastAPI), Java (Spring Boot), TypeScript e PostgreSQL. Tesi di ricerca sull'applicazione dei Large Language Models all'elaborazione del linguaggio naturale e alla traduzione Text-to-SQL. Sono alla ricerca di un ruolo come Backend Developer o posizioni simili.}
\vspace{2pt}


%-----------ISTRUZIONE-----------
\section{Istruzione}
  \resumeSubHeadingListStart
    \resumeSubheading
        {Università di Verona}{Verona, Italia}
        {Laurea Magistrale in Ingegneria e Scienze Informatiche (105/110)}{Ott. 2022 -- Ott. 2025}

        \resumeItemListStart
          \resumeItem{\textbf{Tesi:} \emph{Utilizzo degli LLM per la traduzione del linguaggio naturale in query SQL} --- valutazione e confronto di diversi LLM (SQLCoder, GPT, Gemini, Claude) sulla loro capacità di generare query SQL accurate a partire dal linguaggio naturale, analizzando i compromessi tra accuratezza, latenza e costo tra le diverse famiglie di modelli.}
        \resumeItemListEnd
    \resumeSubheading
      {Università di Verona}{Verona, Italia}
      {Laurea Triennale in Informatica}{Ott. 2019 -- Ott. 2022}

      \resumeItemListStart
          \resumeItem{\textbf{Tesi:} \emph{Sistema di monitoraggio della salute e di suggerimento di attività basato sui dati Fitbit e sull'algoritmo Apriori, sviluppato con il framework Ionic Vue.}}
      \resumeItemListEnd

  \resumeSubHeadingListEnd


%-----------ESPERIENZA-----------
\section{Esperienza}
  \resumeSubHeadingListStart

    \resumeSubheading
      {IT Support Specialist}{Set. 2024 -- Ott. 2025}
      {Generali Italia: Banca e Assicurazioni \normalfont\textit{(Tempo pieno, in parallelo agli studi magistrali)}}{Verona, Italia}
      \resumeItemListStart
        \resumeItem{Fornita assistenza tecnica ai dipendenti tramite la risoluzione di problemi hardware, software e di rete, garantendo tempi di inattività minimi e la regolarità delle operazioni quotidiane.}
      \resumeItemListEnd

    \resumeSubheading
      {Software Developer -- Tirocinio}{Gen. 2024 -- Mag. 2024}
      {BSQ Security}{Verona, Italia}
      \resumeItemListStart
        \resumeItem{Sviluppata un'API RESTful con Flask per un sistema di rilevamento allarmi CCTV basato su intelligenza artificiale, abilitando un'integrazione fluida con client front-end e servizi esterni.}
        \resumeItem{Sviluppato un bot Telegram per inviare notifiche automatiche di allarme agli utenti iscritti, integrato con l'API backend.}
        \resumeItem{Eseguito debugging e unit testing dell'applicazione web qboard.app utilizzando pytest, garantendo l'affidabilità funzionale prima del rilascio.}
        \resumeItem{Redatta documentazione tecnica per assicurare l'affidabilità e la manutenibilità dei componenti sviluppati.}
      \resumeItemListEnd

    \resumeSubheading
      {Vue Developer -- Tirocinio}{Dic. 2022 -- Apr. 2023}
      {RealT}{Verona, Italia}
      \resumeItemListStart
        \resumeItem{Sviluppata un'applicazione web utilizzando Vue.js, TypeScript, Tailwind CSS e Docker.}
      \resumeItemListEnd

  \resumeSubHeadingListEnd


%-----------COMPETENZE TECNICHE-----------
\section{Competenze Tecniche}
 \begin{itemize}[leftmargin=0.15in, label={}]
    \small{\item{
     \textbf{Linguaggi di Programmazione}{: Java, Python, C, SQL (PostgreSQL), JavaScript, TypeScript} \\
     \textbf{Framework}{: Spring Boot, Django, Flask, FastAPI, Node.js, JUnit, Angular, Vue.js, TailwindCSS} \\
     \textbf{Strumenti di Sviluppo}{: Git, Docker, VS Code, IntelliJ, Linux, Postman} \\
     \textbf{Librerie}{: pandas, NumPy, Matplotlib, Transformers, PyTorch, Scikit-learn} \\
     \textbf{Lingue}{: Albanese (Madrelingua), Inglese (Fluente) C2, Italiano (Fluente) C2}
    }}
 \end{itemize}


%-----------PROGETTI-----------
\section{Progetti Personali}
    \resumeSubHeadingListStart
      \resumeProjectHeading
          {\textbf{Sistema di Gestione Inventario per Negozio} $|$ \emph{Python, Django, Angular, PostgreSQL, Docker}}{Mar. 2025 -- Presente}
          \resumeItemListStart
            \resumeItem{Sviluppata un'applicazione web full-stack con backend Django, database PostgreSQL e frontend Angular per gestire inventario, vendite e rifornimenti di un negozio di pescheria al dettaglio.}
            \resumeItem{Distribuita e ospitata l'applicazione tramite container Docker; attualmente in uso attivo da parte del titolare del negozio.}
          \resumeItemListEnd
      \resumeProjectHeading
          {\textbf{Hospitium} $|$ \emph{Java, Spring Boot, Gradle, H2}}{Lug. 2024}
          \resumeItemListStart
            \resumeItem{Sviluppato un sistema software per la gestione delle operazioni di una clinica medica, a supporto di medici e infermieri nella gestione di farmaci e flussi di analisi del sangue.}
            \resumeItem{Implementata un'interfaccia segretariale condivisa per la gestione di appuntamenti e cartelle cliniche dei pazienti.}
          \resumeItemListEnd
      \resumeProjectHeading
          {\textbf{Vue Job Search} $|$ \emph{TypeScript, Vue.js, Django, PostgreSQL, REST API}}{Set. 2022}
          \resumeItemListStart
            \resumeItem{Sviluppata una piattaforma di ricerca lavoro che permette agli utenti di filtrare gli annunci per località, tipo di lavoro e parole chiave.}
            \resumeItem{Implementate funzionalità di autenticazione utente e gestione del profilo.}
          \resumeItemListEnd
      \resumeProjectHeading
          {\textbf{Altri progetti:} \href{https://github.com/LedjoLleshaj}{\underline{github.com/LedjoLleshaj}}}{}
    \resumeSubHeadingListEnd


%-----------RICONOSCIMENTI & FORMAZIONE-----------
\section{Riconoscimenti \& Formazione}
  \resumeSubHeadingListStart
    \resumeSubheading
      {CyberChallenge.IT}{Feb. 2022 -- Giu. 2022}
      {\href{https://cyberchallenge.it/}{\underline{cyberchallenge.it}}}{Università di Verona, Italia}
      \resumeItemListStart
        \resumeItem{Selezionato tramite una prova competitiva nazionale per partecipare a un programma di formazione in cybersecurity di 4 mesi, comprendente Sicurezza delle Reti, Sicurezza delle Applicazioni Web, Crittografia ed Ethical Hacking.}
      \resumeItemListEnd
  \resumeSubHeadingListEnd


%-------------------------------------------
\end{document}
